% Source: https://artofproblemsolving.com/community/c3h1644305_help_to_slow
% Stable identity: aops-topic-1644305-q01
Put $t=\dfrac{x}{x-4}$ (for $x\ne4$). Then $x$ is a root of the second equation exactly when $t$ is a root of $at^2+bt+c=0$, so $t\in(3,8)$. Solving for $x$:
\[
t(x-4)=x\iff x(t-1)=4t\iff x=\dfrac{4t}{t-1}=4+\dfrac{4}{t-1}.
\]
As $t$ runs over $(3,8)$, $t-1$ runs over $(2,7)$, so $\dfrac{4}{t-1}$ runs over $\left(\dfrac{4}{7},2\right)$ and $x$ over $\left(4+\dfrac{4}{7},\,6\right)$. The only integer in this interval is $5$.

It really occurs: $x=5$ corresponds to $t=\dfrac{5}{1}=5\in(3,8)$; for example $a=1$, $b=-9$, $c=20$ (roots $4$ and $5$ of $t^2-9t+20$) give the equation with roots $x=5$ and $x=\dfrac{16}{3}$. The answer is $\boxed{5}$.
